\chapter{Signal Combination}\label{ch:rs:signal-combination}

A research team has forty signals, each with an \omterm{def:rs:anatomy-of-a-predictor:ic}{information coefficient} near 0.015. Blended with equal weights they reach 0.064.
Weighted by their past ICs they reach 0.099 in one simulated market and 0.075 in another; with the textbook weights that maximise the
blend's information ratio, 0.064 and 0.041. In a universe of a hundred names with a year of history, regression weights earn almost nothing in one of the two
markets, and there nothing beats equal weights. Forecasters know this pattern as a puzzle: Stock and Watson (2004) found that the most
successful combinations of output-growth forecasts, like the mean, were the least sensitive to the recent performance of the individual
forecasts. This chapter builds the blends (\texttt{firm.combine}), explains what each assumes, and measures on planted data when the
estimated weights are worth their estimation error.

\section{The geometry of a linear blend}

\begin{definition}[Signal combination, composite signal, equal-weight blend]\label{def:rs:signal-combination:blend}
\emph{Signal combination}\index{signal combination} is the construction of one \omterm{def:rs:anatomy-of-a-predictor:predictor}{predictor} from several. A \emph{composite
signal}\index{composite signal} is a weighted sum $F = \sum_k w_k z_k$ of the signals $z_k$, each standardised across names on every date;
the \emph{equal-weight blend}\index{equal-weight blend} sets every $w_k = 1/K$.
\end{definition}

\begin{proposition}[The IC of an equal-weight blend]\label{prop:rs:signal-combination:geometry}
If $K$ standardised signals each have correlation $\rho$ with the target and pairwise correlation $c$ with each other, their \omterm{def:rs:signal-combination:blend}{equal-weight blend}
has correlation $K\rho/\sqrt{K + K(K - 1)c}$ with the target, which rises with $K$ toward $\rho/\sqrt c$.
\end{proposition}

\begin{proof}
The covariance of $\sum_k z_k$ with the (standardised) target is $K\rho$; its variance is $K + K(K-1)c$. Divide, and let $K \to \infty$.
\end{proof}

The ceiling is set by the correlation between signals, not by their number: forty signals that repeat each other are one signal. The
chapter's first market (\texttt{rs\_combine}, the \emph{stable} world) plants the structure: each month's target is a name's next-month
return on \texttt{firm.synthmkt} (market-adjusted, standardised across the 704 names listed throughout) plus an expected return from eight
latent themes, five of which carry information (from 0.10 down to 0.03) and three none; each of the forty signals measures one theme, five
signals a theme, through a noise its theme shares and a noise of its own whose size grows from 1 to 4 within the theme. Over the last five
years the signals' mean IC is 0.015 and their mean pairwise correlation 0.027. The measured blend of $K$ signals taken in random order follows
the proposition closely (\cref{fig:rs:signal-combination:geometry}): 0.021 for two, 0.042 for ten, 0.064 for all forty against 0.065 from
the formula, on the way to a ceiling of 0.090.

\begin{omfigure}
\begin{tikzpicture}
\begin{axis}[width=10cm,height=5.4cm,xlabel={number of signals blended, $K$},ylabel={out-of-sample IC},tick label style={font=\small},
  label style={font=\small},xmin=0,xmax=41,ymin=0,ymax=0.1,grid=major,grid style={gray!25},
  yticklabel style={/pgf/number format/fixed,/pgf/number format/precision=2},scaled y ticks=false,legend pos=south east,
  legend style={font=\footnotesize},table/col sep=comma]
\addplot[omDef,very thick,mark=*,mark size=1.6pt] table[x=k,y=measured]{figdata/research/14-signal-combination/geometry.csv};
\addplot[omThm,thick,dashed] table[x=k,y=theory]{figdata/research/14-signal-combination/geometry.csv};
\addplot[black!50,thin,dotted] table[x=k,y=ceiling]{figdata/research/14-signal-combination/geometry.csv};
\legend{measured (100 random orders),$K\bar\rho/\sqrt{K + K(K-1)\bar c}$,ceiling $\bar\rho/\sqrt{\bar c}$}
\end{axis}
\end{tikzpicture}

\omcaption{Out-of-sample IC of the \omterm{def:rs:signal-combination:blend}{equal-weight blend} of $K$ of the forty planted signals (stable world, last five years), against the
proposition with their mean IC 0.015 and mean pairwise correlation 0.027. Data: \texttt{rs\_combine} on
\texttt{firm.synthmkt}.}\label{fig:rs:signal-combination:geometry}
\end{omfigure}

\section{Weights from estimates}

\begin{definition}[IC-weighted blend, maximum-ICIR blend]\label{def:rs:signal-combination:weights}
An \emph{IC-weighted blend}\index{IC-weighted blend} sets each weight proportional to the signal's past mean IC (zero if negative). A
\emph{maximum-ICIR blend}\index{maximum-ICIR blend} sets $w \propto \Sigma^{-1}\mu$, with $\mu$ the vector of the signals' past mean ICs and
$\Sigma$ the covariance of their IC time series: the weights that maximise the blended IC's mean over its standard deviation.
\end{definition}

Equal weights assume nothing about the signals; IC weights assume their quality differs and persists; the maximum-ICIR rule assumes, in
addition, that the covariance of forty IC series estimated from 59 months is accurate enough to invert. The second market (the
\emph{drifting} world) removes the first two premises: every signal has the same noise, and each theme's information wanders around a
common mean of 0.04 as an autoregression with a half-life of about two years, so equal weights are right on average and any estimated weights
chase noise and stale information. Fitted on the first five years and judged on the last five, with all 704 names:

\begin{center}
\small
\begin{tabular}{@{}lcccc@{}}
\toprule
 & \multicolumn{2}{c}{stable world} & \multicolumn{2}{c}{drifting world}\\
blend & in sample & out of sample & in sample & out of sample\\
\midrule
equal weights & 0.065 & 0.064 & 0.055 & 0.060\\
IC weights & 0.091 & 0.099 & 0.073 & 0.075\\
maximum-ICIR, no shrinkage & 0.050 & 0.064 & 0.052 & 0.041\\
maximum-ICIR, diagonal $\Sigma$ & 0.090 & 0.098 & 0.074 & 0.074\\
pooled regression (OLS) & 0.093 & 0.100 & 0.076 & 0.071\\
lasso & 0.093 & 0.102 & 0.075 & 0.071\\
the truth (oracle weights) & --- & 0.102 & --- & 0.087\\
\bottomrule
\end{tabular}
\end{center}

With 704 names a month, IC weights capture most of what the truth offers in the stable world, and still beat equal weights in the drifting
one (information that drifts over two years is still there a year later). The unshrunk maximum-ICIR rule is the worst estimated blend in both:
it inverts a covariance matrix of forty noisy series and loads on its smallest eigenvalues. Shrinking that covariance to its diagonal, the shrinkage of Book~4 (chapters 14 and 22) taken to its limit, comes within about a thousandth of the IC-weighted result in both worlds (0.098 against 0.099, 0.074 against 0.075).

\section{Regularised combination}

The weights of a pooled regression of the target on the forty signals are estimated from $59 \times 704$ observations, and there the
regression is as good as IC weights. The puzzle appears when the sample shrinks. In a universe of 100 names with twelve months of history
(\cref{fig:rs:signal-combination:shrink}), unpenalised regression weights earn an out-of-sample IC of 0.073 in the stable world and 0.005 in
the drifting one, against 0.069 for equal weights in both. Shrinking the regression weights toward equal weights,
$w_\lambda = (1 - \lambda)\hat w + \lambda/K$, is the simplest regularisation, and $\lambda$ has an optimum: 0.4 in the stable world (IC 0.081),
1 in the drifting world (equal weights). Smith and Wallis (2009) located the puzzle exactly there: the finite-sample error in estimating the
combining weights, which can exceed the gain from estimating them at all.

\begin{omfigure}
\begin{tikzpicture}
\begin{axis}[width=10cm,height=5.6cm,xlabel={shrinkage toward equal weights, $\lambda$},ylabel={out-of-sample IC},tick label style={font=\small},
  label style={font=\small},xmin=0,xmax=1,ymin=0,ymax=0.1,grid=major,grid style={gray!25},
  yticklabel style={/pgf/number format/fixed,/pgf/number format/precision=2},scaled y ticks=false,legend columns=2,
  legend style={font=\footnotesize,at={(0.5,-0.25)},anchor=north,draw=none,fill=none,
  /tikz/every even column/.append style={column sep=8pt}},table/col sep=comma]
\addplot[omDef,very thick,mark=*,mark size=1.5pt] table[x=lam,y=stable]{figdata/research/14-signal-combination/shrink.csv};
\addplot[omThm,very thick,mark=square*,mark size=1.5pt] table[x=lam,y=drifting]{figdata/research/14-signal-combination/shrink.csv};
\addplot[omDef,thin,dashed] table[x=lam,y=stable_icw]{figdata/research/14-signal-combination/shrink.csv};
\addplot[omThm,thin,dashed] table[x=lam,y=drifting_icw]{figdata/research/14-signal-combination/shrink.csv};
\legend{stable world,drifting world,IC weights (stable),IC weights (drifting)}
\end{axis}
\end{tikzpicture}

\omcaption{A universe of 100 names, blends fitted on twelve months and judged on the next sixty: out-of-sample IC of the pooled regression
weights shrunk toward equal weights ($\lambda = 1$), and of IC weights (dashed). The stable world's optimum is interior ($\lambda = 0.4$); the
drifting world's is equal weights. Data: \texttt{rs\_combine.small\_universe}.}\label{fig:rs:signal-combination:shrink}
\end{omfigure}

\begin{definition}[Forecast combination puzzle]\label{def:rs:signal-combination:puzzle}
The \emph{forecast combination puzzle}\index{forecast combination puzzle} is the repeated empirical finding that simple combinations of
forecasts (the mean) outperform combinations with estimated weights.
\end{definition}

The window matters less than the world: at 6, 12, 24 and 48 months of history the stable world's IC weights beat equal weights (0.085, 0.086,
0.082 and 0.092 against 0.069), and the drifting world's lose until four years of history (0.043, 0.049, 0.067, then 0.074). Ridge and lasso
penalise the regression toward zero rather than toward equal weights; on a standardised panel that shrinks the composite's scale more than its
direction, and in the full panel they moved the out-of-sample IC by at most 0.003 (ridge with a penalty of 1 in the drifting world). The penalty that matters is toward the right prior,
and for signals of similar quality and unstable information the right prior is equal weights.

\section{Stacking}

\begin{definition}[Stacking]\label{def:rs:signal-combination:stacking}
\emph{Stacking}\index{stacking} combines several \omterm{def:rs:anatomy-of-a-predictor:predictor}{predictors} (here, several blends) with weights fitted on their predictions for data not used to
build them: Breiman's (1996) stacked regressions use cross-validated predictions and least squares under non-negativity constraints, an idea
he credits to Wolpert (1992).
\end{definition}

Time-ordered folds are compulsory: a blend fitted on months that follow the months it predicts is a leak. \texttt{firm.stack} fits each base
blend (equal weights, IC weights, regression weights) on an expanding window, collects its predictions for the next block of months, and
regresses the target on them without negative weights. On the full panel the stack put all its weight on IC weights in both worlds, matching
their out-of-sample IC (0.099 and 0.075): \omterm{def:rs:signal-combination:stacking}{stacking} chooses among blends, and it is only as good as the best of them and the length of its
folds.

\section{Orthogonalisation}

\begin{definition}[Orthogonalisation, symmetric orthogonalisation]\label{def:rs:signal-combination:orth}
\emph{Orthogonalisation}\index{orthogonalisation} transforms a set of signals, on each date, into uncorrelated ones spanning the same space:
\emph{sequentially} (Gram--Schmidt), each signal residualised on those before it, so that the order decides who keeps the shared part; or
by \emph{symmetric orthogonalisation}\index{symmetric orthogonalisation}, $Z S^{-1/2}$ with $S$ the signals' correlation matrix, the orthonormal set
closest to the original signals, a construction that goes back to Löwdin's orthonormalisation of overlapping atomic orbitals (1950).
\end{definition}

Orthogonalised signals make weights readable (each weight is the marginal contribution of information no other signal carries) and make
risk-model exposures separable (residualising on a book's factor exposures is the same operation). They do not create information. Equal
weights on the symmetrically orthogonalised signals, recomputed every month, give an out-of-sample IC of 0.064 in the stable world, the same as
the raw signals (0.064), with each orthogonalised signal correlated 0.95 with its original; sequential \omterm{def:rs:signal-combination:orth}{orthogonalisation} in the order of the
in-sample ICs gives 0.050, because the best signal keeps its theme's shared part and the others are reduced to their noise before being
weighted equally with it.

\section{Predictor cards}

\begin{predictorcard}{Composite of forty signals}\label{pred:rs:signal-combination:composite}
\sfield{Definition} $\sum_k w_k z_k$ with $z_k$ standardised by date and $w$ the regression weights shrunk toward equal weights by $\lambda$ chosen on
time-ordered folds.
\sfield{Inputs and timestamps} Each signal as known at the decision date; the weights refitted only on dates before it.
\sfield{Rationale} The equal-weight ceiling $\bar\rho/\sqrt{\bar c}$ (\cref{prop:rs:signal-combination:geometry}); weights earn their keep only when
quality differs persistently.
\sfield{Horizon and half-life} Those of the constituents, weighted.
\sfield{Normalisation} Standardise the composite by date; neutralise to the book (chapter~6).
\sfield{Failure modes} Unshrunk inverse-covariance weights; refitting with look-ahead; a composite dominated by one theme's repeated signals.
\sfield{Sources} Stock and Watson (2004); Smith and Wallis (2009); \texttt{rs\_combine}.
\end{predictorcard}

\section{Tutorial: forty weak signals}\label{tut:rs:signal-combination:tut}

\begin{tutorial}
\textbf{Goal.} Blend forty planted signals in two worlds with every rule of the chapter, in sample and out, on the full panel and on a small
universe, and find the shrinkage toward equal weights that wins. \textbf{End state:}
\cref{fig:rs:signal-combination:geometry,fig:rs:signal-combination:shrink}; the table.
\begin{tutsteps}
\item \textbf{The maximum-ICIR weights}, with the covariance shrunk toward its diagonal.
\omcode{code/firm/combine/firm_combine.py}{68}{75}{Maximum-ICIR weights with a shrunk IC covariance.}\label{lst:rs:signal-combination:icir}
\item \textbf{\omterm{def:rs:signal-combination:orth}{Symmetric orthogonalisation}} of one date's cross-section.
\omcode{code/firm/combine/firm_combine.py}{157}{161}{The orthonormal signals closest to the originals.}\label{lst:rs:signal-combination:sym}
\item \textbf{Run} \texttt{results()} for both worlds, \texttt{small\_universe()} for windows of 6 to 48 months, \texttt{orthogonalised()},
\texttt{geometry()} and \texttt{fig\_combine.py}.
\end{tutsteps}
\textbf{What to change next.} Make the drifting world's half-life six months and watch IC weights lose even on the full panel; give three
signals a negative IC and compare IC weights (which drop them) with equal weights (which keep them).
\end{tutorial}

\section{Build: the blender}\label{bld:rs:signal-combination:combine}

\begin{build}
\bfield{Purpose} One place where the firm's signals become a composite, with every rule of the chapter behind the same interface, and the
out-of-sample comparison that chooses among them.
\bfield{Interface} \texttt{zscore}, \texttt{ic\_series}, \texttt{ic\_matrix}, \texttt{equal}, \texttt{blend(X, w)}, \texttt{ic\_weights},
\texttt{max\_icir(ic, shrink)}, \texttt{toward\_equal(w, lam)}, \texttt{ridge}, \texttt{lasso}, \texttt{nnls}, \texttt{stack(preds, y, folds)},
\texttt{time\_folds}, \texttt{orth\_sequential}, \texttt{orth\_symmetric}, \texttt{residualise\_on}.
\bfield{Rules} Signals are standardised by date before any weight is applied; weights and hyperparameters are fitted on dates before the ones they
are judged on; the \omterm{def:rs:signal-combination:blend}{equal-weight blend} is always reported alongside.
\bfield{Acceptance tests} \texttt{code/firm/combine/tests/}: z-scores and ICs by hand; a planted panel where the good signal earns the largest IC weight,
the diagonal maximum-ICIR equals mean over variance, ridge shrinks, lasso zeroes the noise signal; NNLS by hand; a stack that picks the informative
base; both orthogonalisers orthonormal, the symmetric one closest to the originals, the sequential one keeping its first column.
\bfield{Stretch} Weights by regime; hierarchical blends (within themes, then across); Bayesian weights with a prior at equal weights.
\end{build}

\begin{omsources}
\item J.~H.~Stock and M.~W.~Watson, ``Combination forecasts of output growth in a seven-country data set'', \emph{Journal of Forecasting}
23(6), 2004.
\item J.~Smith and K.~F.~Wallis, ``A simple explanation of the forecast combination puzzle'', \emph{Oxford Bulletin of Economics and Statistics}
71(3), 2009.
\item L.~Breiman, ``Stacked regressions'', \emph{Machine Learning} 24, 1996; D.~H.~Wolpert, ``Stacked generalization'', \emph{Neural Networks}
5(2), 1992.
\item P.-O.~Löwdin, ``On the non-orthogonality problem connected with the use of atomic wave functions in the theory of molecules and
crystals'', \emph{Journal of Chemical Physics} 18(3), 1950.
\end{omsources}

\section{Exercises}

\begin{exercise}[$\star$]\label{exo:rs:signal-combination:1}
Ten signals each have an IC of 0.02 and a pairwise correlation of 0.25. What is the IC of their \omterm{def:rs:signal-combination:blend}{equal-weight blend}, and its limit as the number of
such signals grows?
\end{exercise}

\begin{exercise}[$\star$]\label{exo:rs:signal-combination:2}
Three signals have past mean ICs of 0.03, 0.01 and $-0.01$. Give their IC weights, and their weights shrunk halfway toward equal weights.
\end{exercise}

\begin{exercise}[$\star$]\label{exo:rs:signal-combination:3}
Why must \omterm{def:rs:signal-combination:stacking}{stacking} folds be ordered in time for financial data, when Breiman's \omterm{def:rs:signal-combination:stacking}{stacking} used ordinary cross-validation?
\end{exercise}

\begin{exercise}[$\star\star$]\label{exo:rs:signal-combination:4}
Two signals have ICs $\rho_1 = 0.03$, $\rho_2 = 0.02$ and correlation $c = 0.5$. Find the weights of the best linear blend (for correlation with the
target) and its IC, and compare with equal weights.
\end{exercise}

\begin{exercise}[$\star\star$]\label{exo:rs:signal-combination:5}
Show that sequential \omterm{def:rs:signal-combination:orth}{orthogonalisation} of two standardised signals with correlation $c$ leaves the first unchanged and replaces the second by
$(z_2 - cz_1)/\sqrt{1 - c^2}$. What is the IC of the new second signal in terms of $\rho_1$, $\rho_2$ and $c$?
\end{exercise}

\begin{exercise}[$\star\star$]\label{exo:rs:signal-combination:6}
Why does inverting the covariance of forty IC series estimated from sixty months produce extreme weights? Relate it to the eigenvalues of a sample
covariance matrix (Book~4, chapter~22).
\end{exercise}

\begin{exercise}[$\star\star\star$]\label{exo:rs:signal-combination:7}
\emph{Coding.} With \texttt{small\_universe}, find the optimal $\lambda$ in the stable world for windows of 6, 24 and 48 months. Does the optimum move
toward the regression weights as the window lengthens?
\end{exercise}

\begin{exercise}[$\star\star\star$]\label{exo:rs:signal-combination:8}
\emph{Find the flaw.} ``We fitted the blend weights by regression on the full ten years and they beat equal weights by 40\% over the same period.''
\end{exercise}

\section{Problem: Forty Weak Signals}

\begin{problem}[{Weekend problem --- when weights are worth estimating}]\label{pb:rs:signal-combination:1}
The chapter's two planted worlds on \texttt{firm.synthmkt}: forty signals in eight themes, 704 names, ten years of months split in two; then a
universe of 100 names.

\medskip\noindent\textbf{Part I --- The geometry.}
\begin{enumerate}
\item What are the signals' mean out-of-sample IC and mean pairwise correlation in the stable world?
\item What does \cref{prop:rs:signal-combination:geometry} give for forty signals, and what was measured?
\item What is the ceiling, and why is it not reached by adding signals?
\item Which planted feature makes the signals correlated?
\end{enumerate}

\medskip\noindent\textbf{Part II --- Full panel.}
\begin{enumerate}[resume]
\item Give the out-of-sample ICs of equal weights, IC weights and the truth in both worlds.
\item Why does the unshrunk maximum-ICIR rule do worst?
\item What does shrinking its covariance to the diagonal give?
\item Why do IC weights still beat equal weights in the drifting world?
\end{enumerate}

\medskip\noindent\textbf{Part III --- Small universe.}
\begin{enumerate}[resume]
\item With 100 names and twelve months, give the out-of-sample IC of unpenalised regression weights in each world.
\item State the \emph{named result}: the out-of-sample IC of the best blend against equal weights, and the shrinkage toward equal weights that
maximises it, in each world.
\item How do IC weights fare against equal weights at 6, 12, 24 and 48 months in each world?
\item What did Smith and Wallis identify as the explanation of the puzzle?
\end{enumerate}

\medskip\noindent\textbf{Part IV --- Beyond.}
\begin{enumerate}[resume]
\item What weights did \omterm{def:rs:signal-combination:stacking}{stacking} choose, and why could it do no better?
\item What do equal weights on symmetrically and on sequentially orthogonalised signals give, and why do they differ?
\item When should a team orthogonalise its signals?
\item How would you choose $\lambda$ in production?
\item What would make equal weights the wrong prior?
\item Why do ridge and lasso change the full-panel result so little here?
\item What should every blend report beside its own IC?
\item In one sentence: when are estimated weights worth estimating?
\end{enumerate}
\end{problem}

\section{Interview questions}

\begin{interviewq}[$\star$ \iqroles{researcher}]\label{iq:rs:signal-combination:1}
You have fifty alpha signals. How do you combine them?
\end{interviewq}

\begin{interviewq}[$\star\star$ \iqroles{researcher}]\label{iq:rs:signal-combination:2}
What is the \omterm{def:rs:signal-combination:puzzle}{forecast combination puzzle}, and what explains it?
\end{interviewq}

\begin{interviewq}[$\star\star$ \iqroles{researcher, mle}]\label{iq:rs:signal-combination:3}
How would you use \omterm{def:rs:signal-combination:stacking}{stacking} for alpha signals without leaking?
\end{interviewq}

\begin{interviewq}[$\star\star$ \iqroles{researcher}]\label{iq:rs:signal-combination:4}
Adding a new signal with IC 0.02 did not improve your composite. Why might that be?
\end{interviewq}

\begin{interviewq}[$\star\star$ \iqroles{researcher, risk}]\label{iq:rs:signal-combination:5}
What is the difference between sequential and \omterm{def:rs:signal-combination:orth}{symmetric orthogonalisation}, and when would you use each?
\end{interviewq}

\begin{interviewq}[$\star\star\star$ \iqroles{researcher}]\label{iq:rs:signal-combination:6}
Derive the weights that maximise the correlation of a linear blend with the target, given the signals' ICs and correlation matrix, and explain
why they are dangerous to estimate.
\end{interviewq}
